\chapter*{ВИСНОВКИ}
\addcontentsline{toc}{chapter}{ВИСНОВКИ}

У даній роботі досліджено та реалізовано метод трасування променів у дискретному воксельному просторі з використанням загального світового воксельного об'єму для графічного рендерингу в реальному часі.

\textbf{Основні результати:}

\begin{enumerate}[itemsep=0pt, topsep=0pt]
    \item Розроблено модульну архітектуру системи на базі Bevy (Rust) та WGSL з повним пайплайном відкладеного рендерингу: геометричний прохід для заповнення G-буфера та світловий прохід для м'яких тіней та ambient occlusion.

    \item Реалізовано інверсивне мапування для вокселізації орієнтованих об'єктів та побітову оптимізацію через кодування Мортона (зменшення пам'яті в 8 разів). Впроваджено ієрархічне трасування з MIP-рівнями для швидкого пропуску порожніх регіонів.

    \item Вирішено проблеми числової точності (тіньове та екстремальне тіньове акне) через комбінований bias та впроваджено стохастичні техніки освітлення з uniform disk sampling та cosine-weighted hemisphere sampling.
\end{enumerate}

Запропонований підхід забезпечує простоту реалізації, константний час доступу до вокселів та природну підтримку динамічних сцен. Основні обмеження: втрата інформації про матеріали при побітовому представленні, обмеження розміром сцени та компроміс між точністю, пам'яттю та продуктивністю.

\textbf{Моменти для покращення:}

Поточна реалізація має резерви для оптимізації продуктивності через voxel octree, що заснована на MIP рівнях, але ще не дає максимальної швидкодії. Текстура блакитного шуму завантажується, але не використовується --- її інтеграція покращила б якість стохастичного семплювання разом із темпоральних згладжуванням, зменшуючи артефакти псевдовипадкових генераторів.